\section{Discussion}
\label{sec:discussion}

\subsection{Limited detection by existing validation workflows}
\label{sec:discussion-detection}

The central empirical observation of this study is that the selected validation
workflows detected 2 of 23 confirmed non-equivalent
reproducibility-relevant mutations. The remaining 21 mutations changed the
relevant semantics of the evaluated case but were not rejected by the frozen
validation workflow.

This result should be interpreted as a property of the evaluated
repository--mutation--workflow combinations rather than as a population-wide
estimate for ML research software. The corpus was constructed using a frozen,
operator-specific screening procedure, but the final meaningful denominator is
small and the operator-specific sample sizes are uneven. We therefore do not
infer that a fixed proportion of ML repositories would behave similarly under
other mutations or validation workflows.

Nevertheless, the observed survival of confirmed non-equivalent mutations
shows that successful execution of an existing repository workflow does not
necessarily imply sensitivity to reproducibility-relevant experimental
changes. A workflow may exercise the affected code path while still accepting
a changed seed, split, fold count, or resolved dependency environment.

The result is therefore better understood as evidence about \emph{validation
sensitivity}: in 21 of the 23 confirmed non-equivalent evaluations, the
selected existing workflow did not distinguish the original configuration from
the controlled, semantically different one.


\subsection{Mutation survival is not repository irreproducibility}
\label{sec:discussion-survival}

A SURVIVED mutation has a deliberately narrow meaning in our study. It means
that a confirmed non-equivalent mutation was not detected by the selected
frozen validation workflow. It does not imply that the repository is
irreproducible, that its published results are incorrect, or that no other
repository workflow would detect the same change.

This distinction is especially important because the selected workflows vary in
purpose. Some are upstream tests, whereas others are documented experiments,
examples, validation commands, or CI workflows. A repository may contain
additional checks that were not selected by the frozen protocol, and a
documented example may have been designed to demonstrate usage rather than to
validate numerical invariants.

Conversely, a KILLED mutation should not be interpreted as proof that a
repository is fully protected against reproducibility failures. A KILLED
outcome establishes only that the selected workflow rejected the particular
applied mutation under the study conditions.

Mutation testing is therefore used here as a probe of existing validation
behavior, not as a binary classifier of repository reproducibility.


\subsection{Workflow form and oracle strength}
\label{sec:discussion-oracles}

RQ2 did not reveal a monotonic relationship between the explicitness of the
validation oracle and mutation detection. In the prospectively classified B02
subset, neither of the two meaningful \texttt{assertion} cases was detected,
whereas two of thirteen \texttt{completion-only} cases were detected.

Taken alone, these proportions could be misleading. Both detected B02
mutations occurred in the same
\texttt{dependency-pin} $\times$ \texttt{completion-only} cell. No KILLED
mutation occurred in the other observed operator--oracle combinations.
Moreover, only two of seven selected \texttt{assertion} cases entered the
confirmed non-equivalent meaningful denominator, compared with thirteen of
twenty-two \texttt{completion-only} cases.

These observations prevent attributing the observed RQ2 pattern to oracle
strength. In particular, the study does not show that completion-only workflows
are stronger than assertion-based workflows.

The dependency result also illustrates why workflow completion can still act
as an effective detector for some classes of change. A dependency mutation can
alter the resolved environment in a way that causes imports, APIs, or runtime
behavior to fail before an explicit research-result assertion is reached. In
such a case, a completion-only workflow can kill the mutation even though it
contains no explicit numerical oracle.

This distinction suggests that workflow type, oracle type, and mutation type
capture different properties and should not be collapsed into a single notion
of ``test strength.''


\subsection{Executability as a prerequisite for mutation-based evaluation}
\label{sec:discussion-executability}

A second observation concerns executability. Only 13 of the 39 frozen
cases were evaluable during primary execution. Bounded restoration increased
this number to 24, leaving 15 cases non-evaluable.

These failures are methodologically distinct from mutation outcomes. A
repository that cannot be executed in the contemporary study environment
cannot be assigned a meaningful KILLED or SURVIVED result under the frozen
baseline-first protocol.

At the same time, the scale of the restoration problem is itself relevant to
empirical work on historical ML research software. We encountered barriers
associated with dependency resolution, obsolete package ecosystems, unavailable
or incompatible builds, removed APIs, undeclared runtime requirements, native
dependencies, specialized hardware assumptions, and unavailable workflow
requirements.

We do not interpret these barriers as direct evidence that the corresponding
research results are non-reproducible. They instead demonstrate that
retrospective execution can be constrained by software-ecosystem drift before
the scientific behavior of interest can even be evaluated.

The D027 layer substantially increased the evaluable sample while preserving
the frozen repository revision, candidate mutation, workflow, and oracle
classification. Separating restoration from the canonical primary result was
therefore important: it allowed us to recover additional empirical evidence
without rewriting the original execution record.


\subsection{Implications for ML research-software validation}
\label{sec:discussion-implications}

The results suggest that reproducibility-relevant validation may require checks
that are more specific than successful workflow completion alone. In
particular, workflows intended to guard experimental reproducibility can
benefit from making important experimental assumptions observable and
testable.

Examples include explicit checks on generated partitions, evaluation-protocol
parameters, expected dependency environments, or other artifacts that encode
experiment configuration. The appropriate safeguard depends on the scientific
workflow; the study does not imply that every repository should assert exact
numerical outputs or fix every dependency indefinitely.

A practical role for reproducibility-oriented mutation testing is therefore to
ask a targeted question: if a relevant experimental choice changes, does an
existing validation workflow notice?

This question complements conventional software testing. Traditional unit and
integration tests may establish functional correctness for APIs and components,
while reproducibility-oriented mutations probe whether validation is sensitive
to changes in experimental configuration that may still leave the software
operational.

For research-software developers, surviving confirmed non-equivalent mutations
can identify locations where existing validation does not currently constrain a
reproducibility-relevant choice. Such observations can motivate additional
checks where those choices are scientifically important, without implying that
every surviving mutation represents a defect.


\subsection{Role of MLReproMutate}
\label{sec:discussion-tool}

MLReproMutate was used in this study to make the mutation process explicit,
repeatable, and auditable. The empirical contribution is not merely the
application of mutation testing to ML software, but the use of a
reproducibility-oriented mutation model together with frozen repository
revisions, frozen validation workflows, explicit semantic-equivalence handling,
and structured execution evidence.

This separation between mutation generation, repository execution, semantic
verification, and analysis is important for future extensions. Additional
operators or repository populations can be studied without redefining the
meaning of the outcomes used here, while the current frozen corpus remains an
auditable empirical record.
